\chapter{A real oblivious transfer: the dual-mode OT of Peikert, Vaikuntanathan and Waters}
\label{chap:pvw}

In the protocol of \cite{JKO13} the prover's choice bits are fixed by $\mathcal{F}_{\mathrm{COT}}$
before the garbled circuit is sent, and this is what makes static authenticity enough
(Chapter~\ref{chap:jko}). Jawurek, Kerschbaum and Orlandi describe how to realize
$\mathcal{F}_{\mathrm{COT}}$ with the OT of Peikert, Vaikuntanathan and Waters \cite{PVW08}, adding
an open-all phase in which the sender reveals its coins \cite[Appendix~A]{JKO13}. In this chapter we take
its DDH instance \cite[Sect.~5.3]{PVW08} in messy mode, where the receiver's choice can be
extracted from its first message with a trapdoor and the sender's security is statistical
\cite[Theorem~4.1]{PVW08}. We show that soundness then reduces exactly, with no computational step,
to the end-to-end bound of Theorem~\ref{thm:e2e}. For zero knowledge, the receiver's security is
computational in messy mode; we reduce it to the indistinguishability of the two modes of the
common string, which \cite[Theorem~5.2]{PVW08} derives from DDH.

\section{The dual-mode cryptosystem}

We write the group additively: $\mathbb{G}$ has prime order $q$ and is a $\mathbb{Z}_q$-module.

\begin{definition}[Messy common string, keys, encryption]
  \label{def:pvw}
  \lean{GcZk.PVW.Crs, GcZk.PVW.Crs.Messy, GcZk.PVW.keyGen, GcZk.PVW.enc, GcZk.PVW.dec}
  \leanok
  A common string is a tuple $(g_0, h_0, g_1, h_1)$ of group elements. It is messy, with trapdoor
  $(x_0, x_1)$, if $g_0, g_1 \neq 0$, $h_b = x_b \cdot g_b$ and $x_0 \neq x_1$; only the
  extraction uses the trapdoor. The receiver with choice
  $\sigma$ and secret $r$ sends $\mathit{pk} = (r \cdot g_\sigma, r \cdot h_\sigma)$. The encryption
  of $m \in \mathbb{G}$ on branch $b$ with coins $(s, t)$ is
  \[ (s \cdot g_b + t \cdot h_b,\; s \cdot \mathit{pk}_1 + t \cdot \mathit{pk}_2 + m), \]
  and $(c_0, c_1)$ decrypts with $r$ to $c_1 - r \cdot c_0$. This is \cite[Sect.~5.3]{PVW08},
  where encryption is built on the algorithm $\mathrm{Randomize}$ of \cite[Lemma~5.1]{PVW08}.
\end{definition}

\begin{lemma}[Correctness]
  \label{lem:pvw-dec}
  \lean{GcZk.PVW.dec_enc}
  \leanok
  \uses{def:pvw}
  Decryption is correct on the receiver's branch.
\end{lemma}

\begin{proof}
  \leanok
  $s \cdot (r \cdot g_\sigma) + t \cdot (r \cdot h_\sigma) = r \cdot (s \cdot g_\sigma + t \cdot h_\sigma)$.
\end{proof}

\begin{lemma}[Messy branches]
  \label{lem:pvw-messy}
  \lean{GcZk.PVW.enc_bijective}
  \leanok
  \uses{def:pvw}
  If $g_b \neq 0$ and $\mathit{pk}_2 \neq x_b \cdot \mathit{pk}_1$, then for every message $m$
  the encryption on branch $b$ is a bijection from the coins $\mathbb{Z}_q^2$ onto
  $\mathbb{G}^2$. In particular, with uniform coins the ciphertext is uniform whatever $m$ is. This
  is the second case of \cite[Lemma~5.1]{PVW08}, in exact form.
\end{lemma}

\begin{proof}
  \leanok
  Both sides have $q^2$ elements, so it suffices to show injectivity. Encryption is affine in the
  coins. If the difference $(\delta_s, \delta_t)$ of two coins is mapped to $0$, then
  $(\delta_s + \delta_t x_b) \cdot g_b = 0$, so $\delta_s = -\delta_t x_b$, and then
  $\delta_t \cdot (\mathit{pk}_2 - x_b \cdot \mathit{pk}_1) = 0$, so $\delta_t = 0$.
\end{proof}

\begin{lemma}[Extraction]
  \label{lem:pvw-ext}
  \lean{GcZk.PVW.ext, GcZk.PVW.messy_of_ext, GcZk.PVW.enc_zero_key}
  \leanok
  \uses{def:pvw}
  Let $\mathrm{ext}(\mathit{pk})$ be $0$ if $\mathit{pk}_2 = x_0 \cdot \mathit{pk}_1$ and $1$
  otherwise. If the common string is messy and $\mathit{pk}_1 \neq 0$, the branch
  $1 - \mathrm{ext}(\mathit{pk})$ satisfies the hypothesis of Lemma~\ref{lem:pvw-messy}. On the
  other hand, for $\mathit{pk} = (0, 0)$ the second component of every ciphertext is the message.
\end{lemma}

\begin{proof}
  \leanok
  If $\mathit{pk}_2 = x_0 \cdot \mathit{pk}_1 = x_1 \cdot \mathit{pk}_1$, then
  $(x_0 - x_1) \cdot \mathit{pk}_1 = 0$, which is impossible. The second claim is immediate.
\end{proof}

\begin{remark}
  $\mathrm{ext}$ is the complement of the algorithm $\mathrm{FindMessy}$ of
  \cite[Sect.~5.3]{PVW08}, which outputs the messy branch $1 - \mathrm{ext}(\mathit{pk})$. The argument that $\mathit{pk}_2 = x_0 \cdot
  \mathit{pk}_1$ implies $\mathit{pk}_2 \neq x_1 \cdot \mathit{pk}_1$ needs $\mathit{pk}_1 \neq 0$,
  and the key $(0, 0)$ (the pair $(1, 1)$ in multiplicative notation) has no messy branch. The
  version of August 2008 does not check this. The revision of January 2019 does: its sender outputs
  nothing when $\mathit{pk}_1$ is the identity, its $\mathrm{FindMessy}$ treats that case apart, and
  its receiver chooses $r \neq 0$ \cite[Sect.~5.3]{PVW08}. We follow the revision on the sender's
  side: the verifier rejects any key with $\mathit{pk}_1 = 0$. On the receiver's side the honest
  prover of Definition~\ref{def:pvw-hyb} samples $r$ in all of $\mathbb{Z}_q$, as in
  \cite[App.~A]{JKO13}, so with probability at most $k/q$ it sends a key that the verifier rejects;
  no completeness statement is made for this OT. Lemma~\ref{lem:pvw-dec} holds for every $r$.
\end{remark}

\begin{definition}[Decryption mode]
  \label{def:pvw-dec}
  \lean{GcZk.PVW.Crs.Dec, GcZk.PVW.keyGen_dec}
  \leanok
  \uses{def:pvw}
  A common string is in decryption mode if $h_b = x \cdot g_b$ for a common $x$ and
  $g_1 = y \cdot g_0$ with $y \neq 0$.
  Then the key of the choice $1$ with secret $r$ is the key of the choice $0$ with secret $y r$
  (trapdoor key generation, \cite[Sect.~5.3]{PVW08}).
\end{definition}

\section{Soundness}

\begin{definition}[The protocol with the PVW OT]
  \label{def:pvw-prover}
  \lean{GcZk.PVW.Prover, GcZk.PVW.cts, GcZk.PVW.RealAccept}
  \leanok
  \uses{def:pvw, thm:e2e}
  The common string is messy and fixed. The cheating prover sends keys $\mathit{pk}_j$, one for
  each input; the verifier rejects if some $(\mathit{pk}_j)_1 = 0$. Otherwise it sends, for each
  input $j$ and each branch $b$, the encryption on branch $b$ of $\iota(L_j + b\Delta)$, where
  $\iota$ is any map from labels to $\mathbb{G}$, with uniform coins. The rest is the protocol of
  Theorem~\ref{thm:e2e}, except that after the opening the prover also sees all the coins of the
  OT, as in the open-all phase of \cite[Appendix~A]{JKO13}.
\end{definition}

\begin{lemma}[Sampling the messy ciphertext]
  \label{lem:pvw-coins}
  \lean{GcZk.PVW.coinsOf, GcZk.PVW.coinsEquiv, GcZk.PVW.sum_coins}
  \leanok
  \uses{lem:pvw-messy, lem:pvw-ext, def:pvw-prover}
  Fix the labels. Sampling the coins of both branches is the same as sampling the ciphertexts of
  the messy branches and the coins of the other branches: the coins of a messy branch are the
  preimage of its ciphertext.
\end{lemma}

\begin{proof}
  \leanok
  Lemma~\ref{lem:pvw-messy}, input by input.
\end{proof}

\begin{theorem}[Soundness with the PVW OT in messy mode]
  \label{thm:pvw-sound}
  \lean{GcZk.PVW.xOf, GcZk.PVW.ctsOf, GcZk.PVW.toE2E, GcZk.PVW.cts_coinsOf,
    GcZk.PVW.realAccept_iff, GcZk.PVW.prob_accept_le}
  \leanok
  \uses{lem:pvw-coins, lem:pvw-ext, thm:e2e}
  For an unsatisfiable circuit and a messy common string, with the seed, the hashes and the coins
  of the OT uniform, the verifier accepts with probability at most
  \[ \frac{n_1 + 1}{|V| - 1} + \frac{n_1^2 + n_2 + 1}{|D|} , \]
  the bound of Theorem~\ref{thm:e2e}.
\end{theorem}

\begin{proof}
  \leanok
  Fix the sampled messy ciphertexts and the coins of the other branches. The prover is then a
  prover of Theorem~\ref{thm:e2e} with OT input $\mathrm{ext}(\mathit{pk}_j)$: before the opening it
  computes the ciphertext of the extracted branch from the active label, and after the opening it
  computes the coins of the messy branch from the seed (Lemma~\ref{lem:pvw-coins}). Apply
  Theorem~\ref{thm:e2e} and average.
\end{proof}

\begin{remark}
  The reduction is not efficient: it inverts encryption on the messy branch, which amounts to
  computing discrete logarithms. This does not matter here, since the bound of
  Theorem~\ref{thm:e2e} holds for every prover with the given numbers of queries. Lemma~\ref{lem:pvw-dec}
  is not used in the proof; it is the completeness of the OT.
\end{remark}

\section{Zero knowledge}

\begin{definition}[Verifiers with the PVW OT]
  \label{def:pvw-ver}
  \lean{GcZk.PVW.PVer, GcZk.PVW.keys, GcZk.PVW.toCom, GcZk.PVW.chkP}
  \leanok
  \uses{def:pvw, def:com-games}
  A malicious verifier sees the prover's keys $\mathit{pk}_j = \mathrm{keyGen}(w_j, r_j)$; the rest
  is as in Definition~\ref{def:com-games}, with the ciphertexts of both branches in its first
  message and all the coins of the OT in its opening. The prover decrypts the label of its choice
  and decodes it with a left inverse of $\iota$. Its check after the opening: the material is the
  garbling of the opened seed, and every ciphertext is the encryption of the label of its branch
  with the opened coins. It does not use the witness.
\end{definition}

\begin{lemma}[The check is sound]
  \label{lem:pvw-check}
  \lean{GcZk.PVW.checkSound_pvw}
  \leanok
  \uses{def:pvw-ver, lem:pvw-dec, lem:inputs-below, thm:eval-garble}
  If $C$ reads only the inputs $j < k$ and $C(w) = 1$, then whenever the check passes the honest
  prover's label is the $1$-label of the opened seed.
\end{lemma}

\begin{proof}
  \leanok
  The ciphertext of the prover's branch is an encryption under its own key, so it decrypts to the
  active label (Lemma~\ref{lem:pvw-dec}); apply Lemma~\ref{lem:inputs-below} and
  Theorem~\ref{thm:eval-garble}.
\end{proof}

Since the check recomputes every ciphertext from the seed and the coins, the opening cannot be
equivocated, and no binding argument for the open-all phase \cite[Appendix~A]{JKO13} is needed.

\begin{definition}[Hybrids]
  \label{def:pvw-hyb}
  \lean{GcZk.PVW.realP, GcZk.PVW.hybP, GcZk.PVW.pReal, GcZk.PVW.pHyb, GcZk.PVW.sim_toCom}
  \leanok
  \uses{def:pvw-ver, thm:zk-rom}
  For choice bits $x$, the hybrid $\mathrm{Hyb}_x$ sends the keys of $x$ and runs the simulator of
  Theorem~\ref{thm:zk-rom}. The simulator of the protocol is $\mathrm{Hyb}_0$: it sends the keys of
  the choice $0$ and does not use the witness. As a function of the common string, the probability
  $p_x$ that the verifier outputs $1$ in $\mathrm{Hyb}_x$ is computed without the witness (beyond
  the bits $x$) and without any trapdoor.
\end{definition}

\begin{theorem}[Zero knowledge for a fixed common string]
  \label{thm:pvw-zk}
  \lean{GcZk.PVW.real_hyb, GcZk.PVW.zk_pvw}
  \leanok
  \uses{def:pvw-hyb, lem:pvw-check}
  For every common string, every malicious verifier and every witness of a circuit that reads only
  the inputs $j < k$,
  \[ |p_{\mathrm{real}} - p_0| \leq \frac{q_1 + q_2 + q_3}{|R|} + |p_w - p_0| . \]
\end{theorem}

\begin{proof}
  \leanok
  For fixed keys the verifier is a verifier of Chapter~\ref{chap:commitment}, and
  Lemma~\ref{lem:pvw-check} makes Theorem~\ref{thm:zk-rom} apply: the real execution and
  $\mathrm{Hyb}_w$ differ by at most $(q_1 + q_2 + q_3)/|R|$. Then use the triangle inequality.
\end{proof}

\begin{lemma}[Decryption mode]
  \label{lem:pvw-hyb-dec}
  \lean{GcZk.PVW.pHyb_dec}
  \leanok
  \uses{def:pvw-hyb, def:pvw-dec}
  In decryption mode $p_x$ does not depend on $x$.
\end{lemma}

\begin{proof}
  \leanok
  By Definition~\ref{def:pvw-dec}, the map $r_j \mapsto y r_j$ on the inputs with $x_j = 1$ is a
  bijection of the secrets that turns the keys of $x$ into the keys of $0$, and the simulator does
  not use the secrets otherwise.
\end{proof}

\begin{theorem}[Zero knowledge in messy mode, by switching modes]
  \label{thm:pvw-zk-modes}
  \lean{GcZk.PVW.zk_pvw_modes}
  \leanok
  \uses{thm:pvw-zk, lem:pvw-hyb-dec}
  Let the verifier depend on the common string and make at most $Q$ queries. Averaged over a
  nonempty set $S_m$ of common strings, the real execution and the simulator differ by at most
  \[ \frac{Q}{|R|} + |\mathbb{E}_{S_m} p_w - \mathbb{E}_{S_d} p_w| +
     |\mathbb{E}_{S_d} p_0 - \mathbb{E}_{S_m} p_0| \]
  for every set $S_d$ of decryption-mode strings.
\end{theorem}

\begin{proof}
  \leanok
  The first step of the proof of Theorem~\ref{thm:pvw-zk} gives
  $|p_{\mathrm{real}} - p_w| \leq Q/|R|$ for each common string; averaging over $S_m$,
  $|\mathbb{E}_{S_m} p_{\mathrm{real}} - \mathbb{E}_{S_m} p_w| \leq Q/|R|$. Then use the triangle
  inequality through $\mathbb{E}_{S_d} p_w = \mathbb{E}_{S_d} p_0$ (Lemma~\ref{lem:pvw-hyb-dec}).
  Averaging the bound of Theorem~\ref{thm:pvw-zk} itself would leave the term
  $\mathbb{E}_{S_m} |p_w - p_0|$, which the two differences of averages do not control.
\end{proof}

\begin{remark}
  With $S_m$ and $S_d$ the supports of the messy and the decryption-mode setups, the last two terms
  are the advantages of the distinguishers $p_w$ and $p_0$ between the two modes. The next section
  bounds them by DDH advantages.
\end{remark}

\section{The two modes and DDH}

We follow \cite[Sect.~5.1 and 5.3]{PVW08}, in the 2019 revision, in a group of prime order $q > 2$.
The messy setup samples $g_0, g_1 \neq 0$ and distinct nonzero $x_0, x_1$; the decryption setup
samples $g_0 \neq 0$ and nonzero $x, y$, and outputs $(g_0, x g_0, y g_0, x y g_0)$. The DDH
assumption of \cite[Sect.~5.1]{PVW08} says that for random $g, h \neq 0$ the tuples
$(g, h, a g, a h)$ with $a$ random and $(g, h, a g, b h)$ with $a \neq b$ random are
computationally indistinguishable. For a function $F$ of the tuple with values in $[0, 1]$ we write
$\mathrm{Adv}_{\mathrm{DDH}}(F)$ for the absolute value of the difference of the averages of $F$
over the tuples of the first kind and over those of the second.

\begin{lemma}[The setups as tuples]
  \label{lem:pvw-tuples}
  \lean{GcZk.PVW.mk, GcZk.PVW.parM, GcZk.PVW.parRa, GcZk.PVW.parD, GcZk.PVW.parT,
    GcZk.PVW.parRe, GcZk.PVW.crsM, GcZk.PVW.tupRa, GcZk.PVW.crsD, GcZk.PVW.tupRe,
    GcZk.PVW.crsD_dec, GcZk.PVW.lg, GcZk.PVW.sum_messy, GcZk.PVW.sum_dec}
  \leanok
  \uses{def:pvw, def:pvw-dec}
  The messy setup is the distribution of random tuples conditioned on $a, b \neq 0$, and the
  decryption setup is the distribution of DDH tuples conditioned on $a \neq 0$.
\end{lemma}

\begin{proof}
  \leanok
  With $h = x_0 g_0$, $g_1 = a g_0$ and $b = a x_1 / x_0$, the map
  $(g_0, g_1, x_0, x_1) \mapsto (g_0, h, a, b)$ is a bijection onto the tuples with $g, h \neq 0$
  and $a \neq b$ nonzero; similarly $(g_0, y, x) \mapsto (g_0, x g_0, y)$. This is the computation
  in the proof of \cite[Theorem~5.2]{PVW08}, where it is stated as statistical closeness.
\end{proof}

\begin{theorem}[Indistinguishability of the modes]
  \label{thm:pvw-modes}
  \lean{GcZk.PVW.avg, GcZk.PVW.abs_avg_sub_le, GcZk.PVW.ratio_messy, GcZk.PVW.ratio_dec,
    GcZk.PVW.advDDH, GcZk.PVW.abs_modes_le}
  \leanok
  \uses{lem:pvw-tuples}
  For every function $F$ of the common string with values in $[0, 1]$, the averages of $F$ over
  the messy and the decryption setups differ by at most the DDH advantage of $F$, read as a
  function of the tuple, plus $3/q$.
\end{theorem}

\begin{proof}
  \leanok
  By Lemma~\ref{lem:pvw-tuples}, conditioning removes a proportion $2/q$ of the random tuples and
  $1/q$ of the DDH tuples, and the average of a function with values in $[0, 1]$ over a subset
  moves by at most the proportion removed.
\end{proof}

\begin{theorem}[Zero knowledge in messy mode under DDH]
  \label{thm:pvw-zk-ddh}
  \lean{GcZk.PVW.parM_nonempty, GcZk.PVW.abs_avg_sub_avg_le, GcZk.PVW.pHyb_le_one,
    GcZk.PVW.zk_pvw_ddh}
  \leanok
  \uses{thm:pvw-zk, lem:pvw-hyb-dec, thm:pvw-modes}
  Let the verifier depend on the common string and make at most $Q$ queries. Averaged over the
  messy setup, the real execution and the simulator differ by at most
  \[ \frac{Q}{|R|} + \mathrm{Adv}_{\mathrm{DDH}}(p_w) + \mathrm{Adv}_{\mathrm{DDH}}(p_0) +
     \frac{6}{q} . \]
\end{theorem}

\begin{proof}
  \leanok
  As in Theorem~\ref{thm:pvw-zk-modes}, with Theorem~\ref{thm:pvw-modes} applied to $p_w$ and
  $p_0$.
\end{proof}

\begin{remark}
  The distinguishers $p_w$ and $p_0$ are efficient when the verifier is: they sample keys and run
  the verifier with the simulator, and use no trapdoor. Lean has no notion of efficient
  computation, so this is not checked; under DDH their advantages are negligible.
\end{remark}
